\chapter{Technical Aspects}
\label{ch:technical}

\section{Pipeline and tooling}

End to end, C sources become CIR goldens through a pinned CIR-enabled
Clang invocation; the Lean driver regenerates emitted files from a single
source of truth, typechecks them, diffs goldens, fuzzes emitted code
against native binaries, and runs the rejection suites --- one entry
point, parallel tasks, every failure printed, nonzero exit at the end.
Because the build tool does not track all of the generator's inputs,
drift checks are enforced by diffing rather than by rebuild avoidance.
The driver's roster is the wiring: adding a program means adding it
there, with corpus, gate, proof, golden, fuzz, and rejection stage all
following the same standing convention.

\section{Repository layout}

Translations live next to their proofs in per-family homes, with the
evaluator vocabulary, memory model, transfer chain, validator,
tactics, and user specs as separate libraries composed by one
re-export. Tests mirror the pipeline stages: native C corpus, captured
CIR, golden pins, differential fuzzers, and the driver roster. Anything
finer than this paragraph --- file counts, module lists --- belongs to
Chapter~\ref{ch:progress}, which regenerates it, not to this chapter.

\section{Build and pins}

Reproducibility rests on three pins: the Lean toolchain with its Mathlib
revision, and the CIR-enabled Clang build with its upstream SHA, all
recorded with their smoke tests. The Lean side needs no local Mathlib
build since prebuilt artifacts are fetched; the CIR side is needed only
to re-capture goldens, never to typecheck or verify. Live pin values are
snapshot material for Chapter~\ref{ch:progress}, not prose here.
